% §9.1 From Approximate Sets to Approximate Functions
%
% Cross-reference resolution status:
%
% DEFINED here:
%   \label{sec:atom-to-function}                  -- this section
%
% RESOLVED (chapter 5):
%   \Cref{sec:from-atom-to-set}                   -- chapter 5 §5.1
%   \Cref{ch:approximate-sets}                    -- chapter 5
%
% RESOLVED (chapter 8):
%   \Cref{sec:bsc-capacity}                       -- chapter 8 §8.2
%   \Cref{prop:bsc-capacity}                      -- chapter 8 §8.2 result
%
% RESOLVED (chapter 4):
%   \Cref{sec:kronecker}                          -- chapter 4 §4.4
%   \Cref{thm:kronecker}                          -- chapter 4 §4.4 result
%
% FORWARD REFS within chapter 9 (resolve in Tasks 3-5):
%   \Cref{sec:n-ary-channels}                     -- §9.2 (Task 3)
%   \Cref{sec:n-ary-kronecker}                    -- §9.3 (Task 4)
%   \Cref{sec:n-ary-capacity}                     -- §9.4 (Task 5)
%
% §9.1 opens Part III by recasting Part II's predicate framing as one
% point in a larger family: the codomain of an approximate function need
% not be Boolean. Once $T$ is allowed to be anything finite (or even
% countable), the single $(\fprate, \fnrate)$ pair gives way to a full
% channel matrix $Q$, and the chapter 4 Kronecker theorem and chapter 8
% BSC capacity become special cases of broader results that §9.3 and
% §9.4 prove.

\section{From Approximate Sets to Approximate Functions}
\label{sec:atom-to-function}

\subsection*{Recall: an approximate set is a function with a Boolean codomain}

\Cref{sec:from-atom-to-set} opened Part~II by observing that a classical
set $A \subseteq \Universe$ is nothing more than its membership
predicate, a function $\SetIndicator{A} : \Universe \to \{\top, \bot\}$
that returns $\top$ on the elements of $A$ and $\bot$ elsewhere. The
approximate analogue replaces the exact predicate with a noisy one:
each evaluation $\ASet{A}(x)$ is an instance of the noisy-bit type
$\Bernoulli{\BoolType}$, and the per-element rates collapse, under the
chapter~4 axioms, to a single pair $(\fprate, \fnrate)$. The set as
predicate is the framing that all of \Cref{ch:approximate-sets}
inherits, and the framing pinned the codomain to $\{\top, \bot\}$.
Part~III asks what happens when the codomain is larger.

\subsection*{What if the codomain is bigger?}

The predicate framing is not special to Booleans. Nothing in it relies
on the codomain having exactly two values; what it relies on is that
each evaluation is a draw from a distribution on the codomain, and
that the per-query distribution depends only on which input block the
query falls in. Replacing $\{\top, \bot\}$ with any other type $T$
gives the \emph{approximate function},
\[
  \tilde{f} : \Universe \longrightarrow T,
\]
whose evaluation at $x$ is a noisy $T$-valued report. Each evaluation
$\tilde{f}(x)$ is an instance of the lifted type $\Bernoulli{T}$, the
$T$-valued analogue of the noisy bit of Part~I. The lifting is
mechanical: a $\Bernoulli{T}$ instance has the same triple structure
as a $\Bernoulli{\BoolType}$ instance, except that the output alphabet
has $|T|$ values rather than two, and the per-block rate is a
distribution on $T$ rather than a single error probability.

The framing covers a wide range of computational objects. For $T =
\BoolType$, the codomain is $\{\top, \bot\}$ and the approximate
function is exactly the membership predicate of an approximate set;
Part~II is a special case. For $T = \{1, 2, \ldots, K\}$, the codomain
is a finite enumeration, and $\tilde{f}$ is an n-ary classifier whose
output is one of $K$ labels: the natural object for a tagger whose
inputs are images and whose outputs name a category, or for a router
whose inputs are queries and whose outputs name a shard. For $T =
\mathbb{N}$, the codomain is the natural numbers, and $\tilde{f}$ is a
counter or estimator: the natural object for a sketch that maps a
multi-set to its approximate cardinality, or for a frequency
estimator that maps a stream item to its approximate count. In each
case the mathematical structure is the same: a function from a fixed
input domain to a typed codomain, with the per-evaluation distribution
governed by which block of inputs the query belongs to.

The shift is not just an enlargement of the codomain; it is a change
in how the rates are described. A Boolean codomain admits two rates,
$\fprate$ for false positives and $\fnrate$ for false negatives, and
the pair fully specifies the per-block behavior. A larger codomain
admits more. With $|T| = K$ output values, each input block can map
to any of $K$ outputs, and the per-block behavior is a distribution
on the $K$ outputs rather than a single error probability. The
two-by-two channel matrix of \Cref{sec:from-atom-to-set} becomes a
$K \times L$ matrix where $K$ counts input blocks and $L = |T|$
counts output values, with rows summing to one and the $(i, j)$ entry
giving the probability that the $i$-th block produces the $j$-th
output. The matrix is the rate description for an approximate
function in the same way that the pair $(\fprate, \fnrate)$ was for
an approximate set: it is the complete summary of per-query behavior,
collapsed by the chapter~4 axioms from a potentially much larger
parameter space.

\subsection*{Two natural examples}

\paragraph{A 3-class classifier.}
Take $T = \{\text{cat}, \text{dog}, \text{bird}\}$ and let $\tilde{f}$
be a classifier whose input is an image and whose output is one of
the three labels. The input domain $\Universe$ partitions into three
blocks by true class: the cats, the dogs, and the birds. The per-block
behavior is a distribution on the three labels, recorded as a
$3 \times 3$ confusion matrix:
\[
  Q \;=\;
  \begin{pmatrix}
    Q_{\text{cat}, \text{cat}}  & Q_{\text{cat}, \text{dog}}  & Q_{\text{cat}, \text{bird}} \\
    Q_{\text{dog}, \text{cat}}  & Q_{\text{dog}, \text{dog}}  & Q_{\text{dog}, \text{bird}} \\
    Q_{\text{bird}, \text{cat}} & Q_{\text{bird}, \text{dog}} & Q_{\text{bird}, \text{bird}}
  \end{pmatrix}.
\]
Row $i$ records what the classifier does on inputs whose true class
is $i$. The diagonal entries are the per-class accuracies; the
off-diagonal entries are the confusion rates. A high-quality
classifier has diagonal entries close to $1$; a degenerate one that
always says ``cat'' has its first column equal to one and the rest
zero. The two-rate description of an approximate set generalizes to
the nine-entry description of a 3-class classifier, with six free
parameters once the row-sum constraints are imposed. The structure
that Part~II built for two-by-two membership channels is the
model-order-2 specialization of the structure this chapter builds
for $K \times L$ confusion matrices.

\paragraph{A counter.}
Take $T = \mathbb{N}$ and let $\tilde{f}$ be a counter whose input is
a multi-set query and whose output is an estimate $\hat{n}$ of the
true count $n$. The input domain partitions by true count: the block
$B_n$ consists of all multi-sets whose true count is $n$. The
per-block behavior is a distribution on $\mathbb{N}$, recording the
probability that the counter reports $\hat{n}$ when the truth is $n$.
The channel matrix has countably many rows and countably many
columns, but the structure is the same: row $n$ is a distribution on
$\mathbb{N}$ centered (one hopes) near $n$, with a tail width that
records the counter's precision. Errors are not symmetric: the
counter can both over-count and under-count, and the rate description
records both directions of error as a function of the true count. For
practical sketches (CountMin in chapter~10 and HyperLogLog in
chapter~10's later treatment), the row distributions are
concentrated, with most mass on $\hat{n}$ values close to $n$ and an
analytically tractable tail. The two-rate description of an
approximate set generalizes to a row distribution per block, with the
precision of the counter set by the tail width of the row.

\subsection*{What changes structurally}

The lift from Part~II's predicate framing to Part~III's typed framing
preserves more than it changes. The chapter~4 axioms transfer
unchanged: Axiom~1 still requires that error events at distinct
inputs be mutually independent, and Axiom~2 still requires that the
per-block rate be stable across the block. What changes is the size
of the per-block rate. With a Boolean codomain, the per-block rate is
a single number (a Bernoulli parameter); with a $T$-valued codomain,
the per-block rate is a distribution on $T$. The channel matrix grows
from $2 \times 2$ to $K \times L$ accordingly, and the parameter
count grows from $2$ to $K(L - 1)$, since each of the $K$ rows of an
$L$-column row-stochastic matrix has $L - 1$ free entries.

Two chapter-level results scale with the change.
\Cref{thm:kronecker} of \Cref{sec:kronecker} showed that the joint
distribution of $n$ independent queries on a $\Bernoulli{\BoolType}$
instance factors as the Kronecker product of the $n$ per-query
two-by-two matrices. The same factorization carries over, with the
per-query $K \times L$ matrices in place of the two-by-two ones, and
\Cref{sec:n-ary-kronecker} proves the general form. Similarly,
\Cref{prop:bsc-capacity} of \Cref{sec:bsc-capacity} gave the
closed-form capacity $C = 1 - H(\fprate)$ for the binary symmetric
channel. The closed form is specific to the binary case, but the
framework, capacity as the maximum mutual information over input
distributions, generalizes to arbitrary $K \times L$ channels, and
\Cref{sec:n-ary-capacity} develops the generalization (with closed
forms for the symmetric sub-family, and numerical recipes for the
rest).

\subsection*{Where the chapter goes from here}

The rest of the chapter formalizes the generalization in three
steps. \Cref{sec:n-ary-channels} gives the formal definition of an
n-ary channel as a $K \times L$ row-stochastic matrix, fixes the
notation ($K$ uppercase for model-order, distinguishing it from the
lowercase $k$ used for Bloom-filter hash counts and Miller-Rabin
rounds), and works the constraints out from the chapter~4 axioms.
\Cref{sec:n-ary-kronecker} states and proves the n-ary Kronecker
theorem, the headline result of the chapter, which lifts
\Cref{thm:kronecker} from the binary case to general $K$.
\Cref{sec:n-ary-capacity} lifts \Cref{prop:bsc-capacity} from the
BSC to general $K \times L$ channels, with closed forms where they
exist and pointers to Blahut-Arimoto where they do not. The chapter
closes with a bridge to chapter~10's sketches, each of which is an
approximate function with a different choice of codomain $T$ and a
correspondingly different channel matrix.
